\section{Introduction}
\label{sec:intro}

Memory is the binding constraint in cloud computing.  Cloud providers
price memory at \$3--5 per GiB-month, making it 5--10$\times$ more
expensive per byte than storage~\cite{aws2024pricing}.  Unlike CPU
cycles, which can be time-shared across tenants, physical memory is
exclusively held by its owner.  A virtual machine or container that
allocates 16\,GiB of RAM blocks other tenants from using those
physical pages, regardless of whether the application actively
accesses them.  Studies of production workloads at Google and Microsoft
report that 20--50\% of allocated memory is cold at any given
moment~\cite{maas2020adaptive,amaro2023memory}.

Operating systems have responded with transparent memory compression.
macOS Compressed Memory uses the WKdm algorithm to compress inactive
pages in the VM subsystem~\cite{apple2013compressed}.  Linux zswap
compresses pages headed for swap using LZ4 or
zstd~\cite{jennings2013zswap}.  Windows compresses pages in its
memory manager using the Xpress
algorithm~\cite{microsoft2015compression}.  These mechanisms operate
entirely below the allocator, treating each page as an opaque 4\,KiB
or 16\,KiB byte array.

This placement leaves substantial compression opportunity on the
table.  The allocator possesses structural knowledge that the OS
cannot access: it knows which bytes within a page are live objects
versus freed slots containing stale data; it can route allocations
from the same call site to the same physical pages, creating pages
with structurally homogeneous content; and it tracks per-page access
patterns at object-level granularity rather than relying on
coarse-grained hardware access bits.

\sys is a compression-aware memory allocator that exploits this
structural knowledge.  It interposes on \texttt{malloc} and
\texttt{free} via the alloc8 interposition
framework~\cite{berger2001composing}, requiring no application source
changes.  A background compressor thread periodically scans the heap,
compressing pages that have not been accessed recently.  When the
application later accesses a compressed page, a signal handler
transparently decompresses it.

Three techniques improve compression effectiveness over naive
per-page compression:

\begin{itemize}
\item \textbf{Call-site arena segregation.}  \sys hashes the return
  address of each \texttt{malloc} call to select one of four arenas.
  Allocations from the same call site land in the same arena, placing
  structurally similar objects (e.g., all JSON parse tree nodes, all
  B-tree entries) on the same pages.  Structurally homogeneous pages
  compress 1.5 percentage points better than randomly mixed pages in
  our benchmarks (\S\ref{sec:eval:ablation}).

\item \textbf{Zero-on-free.}  When an object is freed, \sys zeros its
  contents.  Small objects ($\leq$128\,B) are zeroed eagerly in
  \texttt{free()}.  Larger objects are zeroed by the compressor thread
  using non-temporal stores to avoid polluting the CPU cache.
  Replacing stale data with zero runs makes partially-filled pages
  highly compressible (\S\ref{sec:design:zero}).

\item \textbf{Adaptive multi-algorithm compression.}  Newly cold pages
  are compressed with LZ4 (1\,GB/s compression throughput).  Pages
  that remain compressed for 60 seconds are recompressed with zstd at
  level~9, which achieves 2--15 percentage points better compression
  ratios (\S\ref{sec:design:adaptive}).
\end{itemize}

\sys achieves 17--44\% RSS reduction on five real-world applications
with negligible throughput impact:

\begin{itemize}
\item \textbf{Memcached}: 25\% RSS reduction during serve phase with
  0\% throughput loss and 54\% faster cold-key re-access
  (\S\ref{sec:eval:apps}).
\item \textbf{SQLite}: 39\% RSS reduction on an in-memory database
  workload.
\item \textbf{DuckDB}: 30\% peak RSS reduction on TPC-H analytics
  queries.
\item \textbf{JSON processing}: 17\% RSS reduction on a 128\,MB
  parsed document.
\item \textbf{Key-value store}: 33\% RSS reduction with 4\,M
  operations per second.
\end{itemize}

This paper makes the following contributions:

\begin{enumerate}
\item The design and implementation of \sys, a drop-in replacement
  allocator that transparently compresses cold heap pages, including
  call-site arena segregation, adaptive multi-algorithm compression,
  and a signal-based fault handler with per-slot decompression contexts
  (\S\ref{sec:design}).

\item A syscall compatibility layer that interposes on 20 system calls
  and C library I/O functions to prevent kernel EFAULT errors when
  buffer pages are under compression monitoring (\S\ref{sec:design:syscall}).

\item A comprehensive evaluation on five real-world applications
  showing 17--44\% RSS reduction (\S\ref{sec:eval}), and an ablation
  study across 17 configurations that isolates the contribution of
  each optimization (\S\ref{sec:eval:ablation}).
\end{enumerate}
